% =============================================================
% 第6章 テストの実施日と結果 (課題9) 章本体
% reports.tex から \input される．
% =============================================================

\section{はじめに}
本章には，Arduinoオーケストラ開発における
単体テスト・結合テスト・システムテストの実施日と結果を記録する．
テスト項目の定義はマネジメント計画書（\texttt{management\_plan.tex}）のMOP表
（SYS-/INT-/UNIT-）を参照．

\section{テスト実施記録}

\begin{longtable}{p{2cm}p{2cm}p{4.5cm}p{1.5cm}p{4cm}}
\toprule
実施日 & 分類 & テスト項目 & 結果 & 備考 \\
\midrule
\endfirsthead

\toprule
実施日 & 分類 & テスト項目 & 結果 & 備考 \\
6月10日 &  &  &  & まだ性能テストする段階まで達しておらずテストはしていない \\
\midrule
\endhead

% --- 各テストの実施記録を記載 ---
% 書式: YYYY-MM-DD & 単体/結合/システム & MOP番号 + 項目名 & 合格/不合格 & 備考 \\
6月17日 & 単体 & UNIT-4：NOTE\_ON送信タイミング誤差 & 合格 & 100回検出，予測と実際の時間誤差の平均約$0.025\,\mathrm{ms}$（目標$\leq \pm 10\,\mathrm{ms}$）\\
6月17日 & 結合 & INT-2：担当B/C/D結合テスト（Arduino→Processing通信） & 合格 & フォトトランジスタによるLED光検出・拍予測・シリアル通信を経てProcessingでカエルの歌の演奏を確認 \\
6月25日 & 結合 & SYS-3：楽譜進行制御 & 合格 & 2分音符・8分音符を含む楽譜の演奏を確認 \\
6月25日 & 結合 & INT-3：輪唱オフセット制御 & 参考合格 & 輪唱形式の合奏動作を実機確認，楽器間の同期ズレはほぼ見られなかった（数値未計測）\\
6月25日 & 結合 & INT-1：テンポ追従演奏（楽器間タイミング） & 参考 & 輪唱合奏時に楽器間の同期ズレほぼ見られず（$\leq\pm20\,\mathrm{ms}$の数値計測は未実施）\\

\bottomrule
\endfoot
\end{longtable}

\section{合否サマリ（任意）}

% 提出時点での合格率や，未実施項目の一覧を記載する場合はここに．
%
% \begin{itemize}
%   \item 単体テスト：X/Y 合格
%   \item 結合テスト：X/Y 合格
%   \item システムテスト：X/Y 合格
% \end{itemize}
